%%%%%%%%%%%%%%%%%%%%%%%%%%%%%%%%%%%%%%%%%%%%%%%%%%%%%%%%%%%%%%%%%%%%%%
% Lớp 8 | Chương IV. Định lí Thalès | Tính chất đường phân giác của tam giác
% Dạng 8H4?17-1: Tính độ dài đoạn thẳng nhờ tính chất phân giác
%%%%%%%%%%%%%%%%%%%%%%%%%%%%%%%%%%%%%%%%%%%%%%%%%%%%%%%%%%%%%%%%%%%%%%


% nguồn: De_AI_Toan8_cuoi_hoc_ki_i_1 (nhập 21/11/2026)
\begin{bt}%[8H4V17-1]%[ma:84b64d44]
	Cho $\triangle ABC$ có $AB = 6\text{ cm}$, $AC = 9\text{ cm}$, $BC = 10\text{ cm}$. Tia phân giác của góc $A$ cắt cạnh $BC$ tại $D$.
a) Tính độ dài các đoạn thẳng $DB$ và $DC$.
b) Tính tỉ số diện tích $\dfrac{S_{\triangle ABD}}{S_{\triangle ACD}}$.
	\loigiai{
		a) Xét $\triangle ABC$ có $AD$ là đường phân giác của góc $A$ ($D \in BC$). Theo tính chất đường phân giác của tam giác ta có: \\
$\dfrac{DB}{DC} = \dfrac{AB}{AC} = \dfrac{6}{9} = \dfrac{2}{3}$. \\
Áp dụng tính chất dãy tỉ số bằng nhau với $DB + DC = BC = 10\text{ (cm)}$, ta được: \\
$\dfrac{DB}{2} = \dfrac{DC}{3} = \dfrac{DB + DC}{2 + 3} = \dfrac{10}{5} = 2$. \\
Suy ra $DB = 2 \cdot 2 = 4\text{ (cm)}$ và $DC = 3 \cdot 2 = 6\text{ (cm)}$. \\
b) Hai tam giác $ABD$ và $ACD$ có chung đường cao hạ từ đỉnh $A$ xuống cạnh $BC$, nên tỉ số diện tích bằng tỉ số hai cạnh đáy tương ứng: \\
$\dfrac{S_{\triangle ABD}}{S_{\triangle ACD}} = \dfrac{DB}{DC} = \dfrac{2}{3}$.
	}
\end{bt}

% nguồn: De_AI_Toan8_cuoi_hoc_ki_i_1 (nhập 21/11/2026)
\begin{bt}%[8H4V17-1]%[ma:d3049bcb]
	Cho $\triangle ABC$ vuông tại $A$ có $AB = 12\text{ cm}$, $AC = 16\text{ cm}$. Đường phân giác của góc $A$ cắt cạnh huyền $BC$ tại $D$.
a) Tính độ dài cạnh $BC$ và các đoạn thẳng $BD$, $CD$.
b) Kẻ $DE \perp AB$ ($E \in AB$). Tính độ dài đoạn thẳng $DE$.
	\loigiai{
		a) Áp dụng định lí Pythagore cho $\triangle ABC$ vuông tại $A$: \\
$BC^2 = AB^2 + AC^2 = 12^2 + 16^2 = 144 + 256 = 400 \Rightarrow BC = 20\text{ (cm)}$. \\
Vì $AD$ là đường phân giác của góc $A$ trong $\triangle ABC$, ta có: \\
$\dfrac{BD}{CD} = \dfrac{AB}{AC} = \dfrac{12}{16} = \dfrac{3}{4}$. \\
Suy ra $\dfrac{BD}{3} = \dfrac{CD}{4} = \dfrac{BD + CD}{3 + 4} = \dfrac{20}{7}$, nên $BD = \dfrac{60}{7}\text{ (cm)}$ và $CD = \dfrac{80}{7}\text{ (cm)}$. \\
b) Vì $DE \perp AB$ và $AC \perp AB$ nên $DE \parallel AC$. Áp dụng hệ quả định lí Thalès trong $\triangle ABC$: \\
$\dfrac{DE}{AC} = \dfrac{BD}{BC} = \dfrac{3}{7} \Rightarrow DE = \dfrac{3}{7} \cdot 16 = \dfrac{48}{7}\text{ (cm)}$.
	}
\end{bt}

% nguồn: De_AI_Toan8_cuoi_hoc_ki_i_1 (nhập 21/11/2026)
\begin{bt}%[8H4V17-1]%[ma:4856d2bf]
	Cho $\triangle ABC$ có $AB = 15\text{ cm}$, $AC = 20\text{ cm}$, $BC = 25\text{ cm}$. Đường phân giác của góc $B$ cắt cạnh $AC$ tại $D$.
a) Tính độ dài các đoạn thẳng $DA$ và $DC$.
b) Chứng minh $\triangle ABC$ vuông tại $A$, từ đó tính diện tích của $\triangle ABD$ và $\triangle BCD$.
	\loigiai{
		a) Vì $BD$ là đường phân giác của góc $B$ trong $\triangle ABC$, theo tính chất đường phân giác ta có: \\
$\dfrac{DA}{DC} = \dfrac{BA}{BC} = \dfrac{15}{25} = \dfrac{3}{5}$. \\
Suy ra $\dfrac{DA}{3} = \dfrac{DC}{5} = \dfrac{DA + DC}{3 + 5} = \dfrac{20}{8} = 2{,}5$. \\
Do đó $DA = 3 \cdot 2{,}5 = 7{,}5\text{ (cm)}$ và $DC = 5 \cdot 2{,}5 = 12{,}5\text{ (cm)}$. \\
b) Ta có $AB^2 + AC^2 = 15^2 + 20^2 = 225 + 400 = 625 = 25^2 = BC^2$. \\
Theo định lí Pythagore đảo, $\triangle ABC$ vuông tại $A$. \\
Diện tích $\triangle ABC$ là $S_{\triangle ABC} = \dfrac{1}{2} \cdot AB \cdot AC = \dfrac{1}{2} \cdot 15 \cdot 20 = 150\text{ (cm}^2\text{)}$. \\
Diện tích $\triangle ABD$ vuông tại $A$ là $S_{\triangle ABD} = \dfrac{1}{2} \cdot AB \cdot AD = \dfrac{1}{2} \cdot 15 \cdot 7{,}5 = 56{,}25\text{ (cm}^2\text{)}$. \\
Diện tích $\triangle BCD$ là $S_{\triangle BCD} = S_{\triangle ABC} - S_{\triangle ABD} = 150 - 56{,}25 = 93{,}75\text{ (cm}^2\text{)}$.
	}
\end{bt}

% nguồn: De_AI_Toan8_cuoi_hoc_ki_i_1 (nhập 21/11/2026)
\begin{bt}%[8H4V17-1]%[ma:299b6622]
	Cho $\triangle ABC$ cân tại $A$ có $AB = AC = 10\text{ cm}$, $BC = 12\text{ cm}$. Kẻ đường cao $AH$ ($H \in BC$). Đường phân giác của góc $B$ cắt $AH$ tại $I$.
a) Tính độ dài các đoạn thẳng $BH$ và $AH$.
b) Tính độ dài các đoạn thẳng $IH$ và $IA$.
	\loigiai{
		a) Vì $\triangle ABC$ cân tại $A$ nên đường cao $AH$ đồng thời là đường trung tuyến, suy ra $BH = \dfrac{1}{2}BC = \dfrac{1}{2} \cdot 12 = 6\text{ (cm)}$. \\
Áp dụng định lí Pythagore cho $\triangle ABH$ vuông tại $H$: \\
$AH^2 = AB^2 - BH^2 = 10^2 - 6^2 = 100 - 36 = 64 \Rightarrow AH = 8\text{ (cm)}$. \\
b) Xét $\triangle ABH$ có $BI$ là đường phân giác của góc $B$ ($I \in AH$). Theo tính chất đường phân giác ta có: \\
$\dfrac{IH}{IA} = \dfrac{BH}{BA} = \dfrac{6}{10} = \dfrac{3}{5}$. \\
Suy ra $\dfrac{IH}{3} = \dfrac{IA}{5} = \dfrac{IH + IA}{3 + 5} = \dfrac{8}{8} = 1$. \\
Vậy $IH = 3\text{ (cm)}$ và $IA = 5\text{ (cm)}$.
	}
\end{bt}

% nguồn: De_AI_Toan8_cuoi_hoc_ki_i_1 (nhập 21/11/2026)
\begin{bt}%[8H4V17-1]%[ma:c8e70030]
	Cho $\triangle ABC$ có $AB = 8\text{ cm}$, $AC = 12\text{ cm}$, $BC = 15\text{ cm}$. Đường phân giác trong của góc $A$ cắt cạnh $BC$ tại $D$, đường phân giác ngoài của góc $A$ cắt đường thẳng $BC$ tại $E$.
a) Tính độ dài các đoạn thẳng $DB$ và $DC$.
b) Tính độ dài các đoạn thẳng $EB$, $EC$ và $DE$.
	\loigiai{
		a) Vì $AD$ là đường phân giác trong của góc $A$ trong $\triangle ABC$ nên: \\
$\dfrac{DB}{DC} = \dfrac{AB}{AC} = \dfrac{8}{12} = \dfrac{2}{3}$. \\
Áp dụng tính chất dãy tỉ số bằng nhau với $DB + DC = BC = 15\text{ (cm)}$: \\
$\dfrac{DB}{2} = \dfrac{DC}{3} = \dfrac{DB + DC}{2 + 3} = \dfrac{15}{5} = 3$. \\
Suy ra $DB = 6\text{ (cm)}$ và $DC = 9\text{ (cm)}$. \\
b) Vì $AE$ là đường phân giác ngoài của góc $A$ trong $\triangle ABC$ và $AB < AC$ nên $E$ nằm trên tia đối của tia $BC$ (tức $B$ nằm giữa $E$ và $C$), do đó $EC - EB = BC = 15\text{ (cm)}$. \\
Theo tính chất đường phân giác ngoài ta có $\dfrac{EB}{EC} = \dfrac{AB}{AC} = \dfrac{2}{3}$. \\
Suy ra $\dfrac{EB}{2} = \dfrac{EC}{3} = \dfrac{EC - EB}{3 - 2} = \dfrac{15}{1} = 15$. \\
Do đó $EB = 30\text{ (cm)}$, $EC = 45\text{ (cm)}$ và $DE = EB + DB = 30 + 6 = 36\text{ (cm)}$.
	}
\end{bt}
